\newif\ifglobal

%%%%%%%%%%%%%%%%%%%%%%%
%%% Compile to view %%%
%%%%%%%%%%%%%%%%%%%%%%%

%%%%%%%%%%%%%%%%%%%%%%%%%%%%%%%%%%%%%%%%%%%%%%%%%%%
%%% Readme:                                     %%%
%%% https://www.overleaf.com/read/bwdjvfjksvjy/ %%%
%%%%%%%%%%%%%%%%%%%%%%%%%%%%%%%%%%%%%%%%%%%%%%%%%%%

% >>> M?: marks a module setting number or important notice

% >>> M4: Set {./relative-path} to external files for this module <<<
\def \modulefiles {./68-PSDET}
% To add an external file, use:
% (Replace '---' with actual filename)
% '\input{\modulefiles/tables/---.tex}' for tables
% '\includegraphics{\modulefiles/figures/---}' for figures
% '\subfile{\modulefiles/---}' for register files


% >>> M1: This module is in global project? <<<
% 'YES' - uncomment; 'NO' - comment out
\globaltrue

\ifglobal
    \documentclass[main\_\_CN.tex]{subfiles}
   
\else
    % Font "FandolSong-Regular" used by ctex does not have GSUB table.
    % It causes warnings that do not affect compilation and can be ignored.
    \PassOptionsToPackage{quiet}{fontspec}

    \documentclass[openany, 10pt]{book}
   
    % >>> M2: In ./00-shared/preamble-trm-module, also find
    %         settings for visibility of todo notes, watermark <<< 
    \usepackage{preamble-trm-module}

    % Enable Chinese typesetting
    \usepackage{ctex}

    % >>> M3: Temporary labels in Chinese
    %         you can add your own labels there <<<
    \myexternaldocument{./00-shared/temp-labels-cn}

\fi %\ifglobal


\setcounter{secnumdepth}{4}


% >>> M5: If this module needs any updates in
% './00-shared/preamble-trm-module.sty'
% (include additional package, etc.), contact your module PM
% or create the file './00-shared/changelog.tex' make the
% necessary updates yourself and document them in this file <<<

% Make text left-aligned and allow latex to break pages naturally, comment out for full justification
\raggedright
\raggedbottom


\begin{document}


% Sets language into which pre-defined bits of text are translated
\selectlanguage{chinese}
\setchineseformatting


% Do not delete this block
\ifglobal
% Add the following front matter if
% this module is outside of global project
\else
    \InputIfFileExists{readme}{}{}
    \clearpage
    \listoftodos
    \listofdonetodos
    \tableofcontents
    %\listoftables
    %\listoffigures
\fi


%%%%%%%%%%%%%%%%%%%%%%%%%
%%% TRM Module Begins %%%
%%%%%%%%%%%%%%%%%%%%%%%%%


\hypertarget{psdet}{}
\chapter{电源检测器 (PSDET)}
\label{mod:psdet}

\section{概述}

\chipname{} 芯片集成了电源检测器，用于监测片上HP系统调压器的电压。具体而言，该功能由一个电压毛刺检测器实现，通过实时监测供电电压，预防异常波动和毛刺攻击。电源检测器可以在HP系统工作期间持续工作。




%%%%%%%%%%%%%%%%%%%%%%%%%%%%%
%%% Feature List from H21 %%%
%%%%%%%%%%%%%%%%%%%%%%%%%%%%%

\section{主要特性}
\label{sec:alias-features}
\begin{itemize}
    \item 在芯片正常工作阶段保持监测，不需要软件干预
    \item 当监测到约 50 ns 以上的电压毛刺后，自动触发系统复位
\end{itemize}


%%%%%%%%%%%%%%%%%%%%%%%%%%%%%%%%%%%%%%
%%%%%%%%%%%%%%%%%%%%%%%%%%%%%%%%%%%%%%

\section{功能描述}

基于电压毛刺检测的电源检测器的结构如下图所示：


\begin{figure}[H]
    \centering
    \includegraphics[width=0.6\textwidth]{\modulefiles/figures/power_detector_p4}
    \caption{\chipname{} 电源检测器结构}
    \label{fig:power-detector-p4}
\end{figure}



\chipname{} 的电源检测器可以在芯片工作周期内保持检查供电网络上的电压，当监测到波动深度大于 100 mV，且持续时间大于 50 ns 的瞬时电压毛刺时，会自动触发系统复位，检测过程无需寄存器配置。由该功能触发的系统复位会被 HP CPU0 复位源\hyperref[fielddesc:LPCLKRSTHPCORE0RESETCAUSE]{LP\_CLKRST\_HPCORE0\_RESET\_CAUSE} 记录，对应 0x1B 复位源。


该功能默认关闭，可通过烧写 \hyperref[fielddesc:EFUSEPVTGLITCHEN]{EFUSE\_PVT\_GLITCH\_EN} 与 \hyperref[fielddesc:EFUSEPVTGLITCHMODE]{EFUSE\_PVT\_GLITCH\_MODE} 开启功能。
更多有关 eFuse 的信息，请参考章节 \ref{mod:efuse} \textit{\nameref{mod:efuse}}。


更多有关复位的信息，请参考章节 \ref{mod:resclk} \textit{\nameref{mod:resclk}}。


\end{document}
